% PROPOSAL ONLY. Requires xcolor. Do not overwrite existing manuscript revisions.
% Suggested insertion after the sample-adaptive reference paragraph in Section3,
% before Figure adaptive_ipnm_reference; alternatively replace only its broad
% "equally or nearly suitable" interpretation after author approval.
\textcolor{green}{The reference is conditional on the specified initialization
and numerical construction, rather than a unique globally optimal benchmark.
Consistent relabeling leaves DI unchanged, but equal preliminary OT losses do
not imply identical final partitions or DI after refinement and hardening.
In a focused check with two prespecified samples of size 8 from the clustered
population ($N=144$, unequal IPs), ten sample-initialized runs gave DI ranges
0.0279 and 0.0043. Alternative initialization, considered only as a sensitivity
stress test, yielded a reference with a 4.45\% native OT-loss gap and 2.65\%
refined soft-compactness gap, changing the second sample's DI from $-0.5822$
to $-0.4220$ without changing the sign of $D_{\mathrm{net}}$. Both losses were
lower for the alternative. These percentages use the stated per-run median
normalization; a common-scale OT comparison instead gives a 20.86\% improvement.
Thus the construction must be reported and fixed for reproducibility; we do
not assert invariance to alternative reference choices.}
